\section{Character classes and the contradiction}\label{sec:characters}

Fix $n\geq1$ and an odd integer $d>n$ with $d\geq3$, as in
\Cref{sec:local-tests}.  We use its data
\[
 W=O\times\R^{d-n}\subset\R^d,\qquad
 c:S^d\longrightarrow W^+,
 \qquad q:W^+\longrightarrow Y=W^+/G,
 \qquad a:Y\longrightarrow S^d
\]
where $O\subset\R^n$ is connected and open, $G=\Zp$ acts faithfully
and preserves orientation on $W$, and $\deg(aqc)=1$.
Put $y_\infty=q(\infty)$.  The orientation of $S^d$ is the one used
for $u_d$ in \Cref{thm:sphere-lattice}.
All sheaves and pairings below are over $\R$; the notation $\C$ denotes
the real coefficient plane with its additional complex scalar operations.
We will associate integral classes to finite partitions of the character
set, show that they are natural under maps and homotopies of tests, and
use the finite quotient sheets to equate their rational images. The final
step compares those images with the fixed lattice from
\Cref{thm:sphere-lattice}.

\subsection{Characters and proper direct images}

Let $\Lambda=\bigcup_{j\geq 0}\mu_{p^j}\subset S^1$ be the countable
group of complex roots of unity of $p$-power order. Define the orbit sheaf
$\mathcal H=q_*\const{\C}{W^+}$. On this sheaf let $T$ denote the
action of $1\in G$, with the convention
$(Ts)(x)=s((-1)x)$, and write $M_\lambda$ for multiplication by
$\lambda\in\C$.  Define
\[
 \mathcal H_\lambda=\ker(T-M_\lambda),\qquad
 \mathcal H_D=\bigoplus_{\lambda\in D}\mathcal H_\lambda
 \quad(D\subset\Lambda).
\]

\begin{lemma}[Character coproducts]\label{lem:character-coproduct}
There are canonical isomorphisms
\begin{equation}\label{char:orbit-sum}
 Rq_*\const{\C}{W^+}\simeq\mathcal H
       \simeq\bigoplus_{\lambda\in\Lambda}\mathcal H_\lambda.
\end{equation}
For every continuous map $t:Y\to X$ to a compact Hausdorff space and
every $D\subset\Lambda$, the natural map
\begin{equation}\label{char:derived-sum}
 \bigoplus_{\lambda\in D}Rt_*\mathcal H_\lambda
       \longrightarrow Rt_*\mathcal H_D
\end{equation}
is an isomorphism in the bounded-below derived category of real sheaves
on $X$.  The sum on the left is a coproduct in the ambient derived
category.
\end{lemma}

\begin{proof}
Apply \Cref{lem:orbit-injectivity} to the compact metrizable
$G$-space $W^+$ with $k=\R$.  The coefficient sheaf
$\const{\C}{W^+}$ is the direct sum of two real constant sheaves,
so additivity gives the first isomorphism in \eqref{char:orbit-sum}.
This argument imposes no dimension bound on $Y$.

The stalk of $\mathcal H$ at an orbit is the complex vector space of
locally constant functions on $G/H$.  The pullback of any such function
to $G$ factors through $G/p^jG$ for some $j$: a finite clopen partition
on which the function is constant admits a refinement by cosets of a
common open subgroup.  The representation is consequently the union
of finite cyclic representations.  Each of the latter decomposes into
the eigenspaces of $T$, since $z^{p^j}-1$ has distinct complex roots.
Kernels and sheaf direct sums can be checked on stalks, proving the
second isomorphism in \eqref{char:orbit-sum}.  This asserts an algebraic
sheaf direct sum; a section over a noncompact open set need not have
globally finite character support.

We record explicitly the compactness argument needed for
\eqref{char:derived-sum}.  If $K$ is compact Hausdorff and
$(\mathcal F_\lambda)$ is a family of sheaves, every section of
$\bigoplus_\lambda\mathcal F_\lambda$ on $K$ involves only finitely many
indices.  Indeed it has this property locally by the definition of
sheafification, and a finite subcover makes the set of indices finite
on $K$.  Thus global sections commute with direct sums.  The same
argument on every closed subset of $K$ shows that direct sums of soft
sheaves are soft: extend the finitely many component sections
separately.  Soft resolutions and their acyclicity on compact Hausdorff
spaces therefore give
\begin{equation}\label{char:compact-cohomology}
 H^r\!\left(K,\bigoplus_\lambda\mathcal F_\lambda\right)
       \simeq\bigoplus_\lambda H^r(K,\mathcal F_\lambda)
 \qquad(r\geq0).
\end{equation}
The soft-resolution facts used here are the compact case of
\cite[Propositions~11.2.4, 11.2.5, and~11.2.8]{SchapiraNotes}.

Apply \eqref{char:compact-cohomology} to $K=t^{-1}(x)$ and the
restrictions of the eigensheaves.  Proper base change identifies the
stalk cohomology of \eqref{char:derived-sum} with these isomorphisms.
All the derived images have the common lower bound zero, so their sum
remains bounded below.  Finally, countable sheaf sums are exact, and
termwise sums compute coproducts in the derived category by
\cite[Lemma~13.33.5]{StacksDerivedSums}.  This proves both assertions.
\end{proof}

For the rest of the construction, a \emph{compact test} means a compact
test space $X$ as in \Cref{def:test-categories}.  We consider a continuous
map $t:Y\to X$ for which
\[
 b_t=t(y_\infty)
\]
has a basis of contractible open neighborhoods in $X$.  Write
\begin{equation}\label{char:test-form}
 f_t=tqc:S^d\longrightarrow X,\qquad
 A_t=Rf_{t*}\const{\C}{S^d},\qquad
 \alpha_t:A_t\xrightarrow{\sim}\D_{X,d}A_t.
\end{equation}
Here $\alpha_t$ is the pushforward of the sphere orientation pairing
with coefficient form
\[
 B(z,w)=\operatorname{Re}(z\overline w).
\]
The object $A_t$ belongs to $\T(X)$, being two copies of a compact
polyhedron generator; its form is supplied by
\Cref{prop:test-duality,prop:test-pushforward}.

Set $X_t^\circ=X\setminus\{b_t\}$.  The preimage
$f_t^{-1}(X_t^\circ)$ lies in $W$, since $c$ sends $S^d\setminus W$ to
$\infty$.  It is precisely $(tq)^{-1}(X_t^\circ)$ under the identity
map on $W$.  The restricted map to $X_t^\circ$ is proper: it is the
base change over an open set of the proper map $f_t$.  Composition and
open base change give a canonical identification
\begin{equation}\label{char:open-identification}
 A_t|_{X_t^\circ}\simeq(Rt_*\mathcal H)|_{X_t^\circ}.
\end{equation}
The pairing on this open set is the direct image of the orientation
pairing on the actual oriented manifold $f_t^{-1}(X_t^\circ)\subset W$.
Thus duality is used on the manifold and on the test space; no duality
on $Y$ or $W^+$ enters this identification.

For $D\subset\Lambda$ put
$B_D=(Rt_*\mathcal H_D)|_{X_t^\circ}$.  Every $B_D$ is a retract of
$A_t|_{X_t^\circ}$ by \eqref{char:orbit-sum} and the complementary
subsum.  Consequently it is bounded with the same cohomological
amplitude bounds as $A_t|_{X_t^\circ}$.  These statements concern the
ambient derived category on $X_t^\circ$; membership of $B_D$ in a test
subcategory on that open set is unnecessary. Their Verdier duals are
bounded as well, since they are retracts of \(\D_{X_t^\circ}(A_t|_{X_t^\circ})\).
The inclusions of bounded and bounded-below derived categories into the
ambient derived category are fully faithful
\cite[Definition~13.11.3 and Lemma~13.11.6]{StacksDerivedCategories}.
Thus the coproduct property in that ambient category detects the
morphisms between bounded objects used below.

\begin{lemma}[Real orthogonality]\label{lem:orthogonality}
If $D,E\subset\Lambda$ are disjoint, the cross pairing between $B_D$
and $B_E$ induced by $\alpha_t|_{X_t^\circ}$ is zero.  In particular,
the projectors associated with any finite partition of $\Lambda$ are
self-adjoint.
\end{lemma}

\begin{proof}
Write $A=A_t|_{X_t^\circ}$ and $\alpha=\alpha_t|_{X_t^\circ}$, and
let $R^\dagger=\alpha^{-1}\D_{X_t^\circ,d}(R)\alpha$ denote the adjoint
of an endomorphism of $A$.  The orientation-preserving action of $G$
on the proper manifold map above acts by isometries of $\alpha$,
by naturality of orientation and proper duality.  Complex scalar
multiplication by a unit also preserves $B$.  In particular, if
$J=M_{\mathrm i}$, then
\[
 J^\dagger=-J,\qquad
 M_\lambda=(\operatorname{Re}\lambda)\id+
               (\operatorname{Im}\lambda)J,\qquad
 T^\dagger=T^{-1},\qquad M_\lambda^\dagger=M_{\lambda^{-1}}
 \quad(\lambda\in\Lambda).
\]

Fix $\lambda\notin E$.  On $B_{\{\lambda\}}$ the endomorphism
$T-M_\lambda$ is zero.  Let $R_E$ denote the endomorphism of $B_E$
induced by $T^{-1}-M_{\lambda^{-1}}$ on $\mathcal H_E$.
The latter operator is invertible: on the $\mu$-summand its
inverse is multiplication by
$(\mu^{-1}-\lambda^{-1})^{-1}$.  These componentwise inverses define
one endomorphism of the algebraic sheaf direct sum, and hence induce
an inverse $R_E^{-1}$ on $B_E$ after $Rt_*$ and restriction.
For the two inclusions $i_\lambda:B_{\{\lambda\}}\to A$ and
$i_E:B_E\to A$, adjointness yields
\[
 \begin{split}
 0
 &=\D_{X_t^\circ,d}(i_E)\alpha(T-M_\lambda)i_\lambda\\
 &=\D_{X_t^\circ,d}(R_E)\,
       \D_{X_t^\circ,d}(i_E)\alpha i_\lambda.
 \end{split}
\]
The first factor in the last line is invertible, so the cross pairing
from $B_{\{\lambda\}}$ to $\D_{X_t^\circ,d}B_E$ vanishes.
This applies also when $\mu=\overline\lambda\ne\lambda$: the common
complex scalar operators distinguish these two real subrepresentations.

By \Cref{lem:character-coproduct} and exact open restriction, $B_D$
is the derived coproduct of the $B_{\{\lambda\}}$, $\lambda\in D$.
A morphism out of a coproduct is determined by its composites with
the coprojections.  The cross pairing from $B_D$ therefore vanishes.
Symmetry gives the cross pairing in the opposite order.  In a finite
decomposition these equalities say precisely that every projector is
self-adjoint.
\end{proof}

\subsection{Integral classes without splitting quotient idempotents}

The projectors constructed above live on the punctured target. We now
use their germs to construct integral classes on the whole target.

Fix a finite partition
$\Lambda=D_0\sqcup\cdots\sqcup D_{m-1}$.
Choose a contractible open neighborhood $U$ of $b_t$, and set
\[
 \mathcal Q_U(X)=\T(X)/\T(U).
\]
The open projectors above define endomorphisms $e_i$ of the image of
$A_t$ in $\mathcal Q_U(X)$ by \Cref{thm:germs}, since
$X_t^\circ$ contains $X\setminus U$.  In this quotient they satisfy
\begin{equation}\label{char:projector-identities}
 e_i e_j=\delta_{ij}e_i,\qquad
 \sum_{i=0}^{m-1}e_i=\id,\qquad e_i^\dagger=e_i.
\end{equation}
Only endomorphisms of the existing object $A_t$ are asserted here.

\begin{definition}[Character classes]\label{def:character-classes}
Let
\[
 \ell_U:E_d(X)\xrightarrow{\sim}W_d(\mathcal Q_U(X))
\]
be the integral isomorphism of \Cref{prop:contractible-quotient}.
For $0\leq i<m$ define $z_i(t)\in E_d(X)$ uniquely by
\begin{equation}\label{char:class-definition}
 \ell_U(z_i(t))=
 [A_t,\alpha_t]+[A_t,\alpha_t\circ(2e_i-\id)].
\end{equation}
\end{definition}

This is a class in the ordinary quotient Witt group: the second form
is nonsingular and symmetric because $2e_i-\id$ is a self-adjoint
involution.  If $e_i$ is represented by an orthogonal direct summand,
the involution is $+\id$ on that summand and $-\id$ on its complement.
The sum in \eqref{char:class-definition} then cancels the complementary
form and doubles the selected form.  An image object for $e_i$ is not
part of the definition.
The definition is independent of $U$.  Indeed, for two choices $U_1,U_2$
choose a contractible neighborhood $U_0\subset U_1\cap U_2$ of $b_t$.
The quotient functors $\mathcal Q_{U_0}(X)\to\mathcal Q_{U_j}(X)$
carry the projectors to each other by their germ descriptions, and
commute with localization from $E_d(X)$.  Since both $\ell_{U_j}$
are injective, the lifts defined using $U_0,U_1,U_2$ coincide.

\begin{proposition}[Sum of the character classes]\label{prop:sum-classes}
Let $t:Y\to X$ be a compact test map as in \eqref{char:test-form},
with a contractible open neighborhood basis at $b_t$. For every finite
partition $\Lambda=D_0\sqcup\cdots\sqcup D_{m-1}$, its character
classes satisfy the integral identity
\begin{equation}\label{char:class-sum}
 \sum_{i=0}^{m-1}z_i(t)=2[A_t,\alpha_t]\quad\text{in }E_d(X).
\end{equation}
\end{proposition}

\begin{proof}
We work in $\mathcal Q_U(X)$ and abbreviate $(A_t,\alpha_t)$ to
$(A,\alpha)$.  On $A^{\oplus m}$ let
\[
 Q=\operatorname{diag}(2e_0-\id,\ldots,2e_{m-1}-\id),
 \qquad Q_0=\operatorname{diag}(\id,-\id,\ldots,-\id).
\]
Let $P_j$ be the permutation matrix interchanging the zeroth and
$j$th copies, with $P_0=\id$, and let $E_j$ be the diagonal matrix
having $e_j$ in every diagonal entry.  Define
\[
 S=\sum_{j=0}^{m-1}E_jP_j.
\]
All the $E_j$ commute with the permutation matrices.  Using only
\eqref{char:projector-identities}, we have $S^2=\id$ and
$S^\dagger=S$, where the adjoint is for $\alpha^{\oplus m}$.
Moreover,
\[
 S^\dagger Q_0 S
  =\sum_{j=0}^{m-1}E_jP_jQ_0P_j
  =Q.
\]
For the last identity, its $i$th diagonal entry is
$e_i-\sum_{j\ne i}e_j=2e_i-\id$ and all off-diagonal entries are zero.
Thus $S$ is an isometry from the sum of the twisted forms to
$(A,\alpha)\perp(A,-\alpha)^{\perp(m-1)}$.  The sum of their Witt
classes is $(2-m)[A,\alpha]$.  Adding the $m$ untwisted terms in
\eqref{char:class-definition} gives $2[A,\alpha]$.  The injectivity
of $\ell_U$ proves \eqref{char:class-sum}.
\end{proof}

\subsection{Covariance and homotopy}

\begin{proposition}[Covariance on compact tests]\label{prop:class-covariance}
Let $t:Y\to X$ be a compact test map as in \eqref{char:test-form},
and let $g:X\to X'$ be a continuous map of compact tests. Suppose $X$
has a contractible neighborhood basis at $b_t$ and $X'$ has one at
$b_{gt}=g(b_t)$.  For every fixed finite partition of $\Lambda$,
\begin{equation}\label{char:covariance}
 g_*z_i(t)=z_i(gt)\quad\text{in }E_d(X').
\end{equation}
\end{proposition}

\begin{proof}
Choose a contractible open $U'\ni g(b_t)$, and then a contractible
open $U\ni b_t$ with $U\subset g^{-1}(U')$.  Pushforward takes an
object supported in a compact subset of $U$ to one supported in a
compact subset of $U'$.  Consequently
\Cref{prop:test-pushforward,thm:witt-localization} give a
duality-preserving quotient functor
\[
 \overline g_*:\mathcal Q_U(X)\longrightarrow\mathcal Q_{U'}(X')
\]
and a commuting square on Witt groups with the maps $\ell_U,\ell_{U'}$.
Composition of proper direct images identifies the pushed total form
with $(A_{gt},\alpha_{gt})$.

We check the endomorphisms in this identification.  Represent a
projector $e_i$ by a quotient roof.  Its restriction agrees with the
character projector on some open neighborhood $N$ of $X\setminus U$.
Shrinking $N$ if necessary, arrange $N\subset X\setminus\{b_t\}$ and
that the denominator of the roof is invertible there.  Define
\begin{equation}\label{char:covariance-neighborhood}
 N'=X'\setminus g(X\setminus N).
\end{equation}
Since $X$ is compact, $g(X\setminus N)$ is closed.  Since
$X\setminus N\subset U$ and $g(U)\subset U'$, the set $N'$ contains
$X'\setminus U'$.  Also $g^{-1}(N')\subset N$, and
$g(b_t)\notin N'$.  Over $N'$ the pushed roof can therefore be
computed entirely where the original projector agrees with its
character description.

Open base change and composition identify that description with the
projector on $R(gt)_*\mathcal H$ for $D_i$.  Every upstairs preimage
in this computation avoids $\infty$, so
\eqref{char:open-identification} identifies it with the projector for
$gt$ on $A_{gt}$.  The germ theorem gives equality of these
endomorphisms in $\mathcal Q_{U'}(X')$.  The quotient functor thus
carries the two forms in \eqref{char:class-definition} to their
counterparts for $gt$.  Applying the integral localization
isomorphisms proves \eqref{char:covariance}.
\end{proof}

\begin{proposition}[Homotopy of character classes]\label{prop:class-homotopy}
For a fixed partition of $\Lambda$, homotopic maps
$t_0,t_1:Y\to S^d$ satisfy $z_i(t_0)=z_i(t_1)$ in $E_d(S^d)$.
\end{proposition}

\begin{proof}
The value at $y_\infty$ may move during the homotopy.  We will combine
two sufficiently close stages into a single compact test and recover
them by homotopic projections.
Consider the compact test
\begin{equation}\label{char:close-pair}
 C_d=\{(x_0,x_1)\in S^d\times S^d:
                     \langle x_0,x_1\rangle\geq0\}.
\end{equation}
It is compact semialgebraic, hence homeomorphic to a finite polyhedron.
It has a contractible open neighborhood basis at each point.  To
check this also at its boundary, the differential of
$F(x_0,x_1)=\langle x_0,x_1\rangle$ on $S^d\times S^d$ is nonzero
where $F=0$: its derivative in the tangent direction $x_0$ at $x_1$
equals $1$.  Thus $C_d$ is a smooth $2d$-manifold with boundary, with local
Euclidean or half-space charts; sufficiently small balls in these
charts give the required basis.  Its test forms have degree $d$,
inherited by pushforward from $S^d$; the target dimension is $2d$.

The projections $\pi_0,\pi_1:C_d\to S^d$ are homotopic by
\[
 (x_0,x_1,r)\longmapsto
 \frac{(1-r)x_0+rx_1}{\|(1-r)x_0+rx_1\|},\qquad 0\leq r\leq1.
\]
The denominator is nonzero, since its square is
$(1-r)^2+r^2+2r(1-r)\langle x_0,x_1\rangle\geq1/2$.

Let $H:Y\times[0,1]\to S^d$ be the given homotopy.  Uniform continuity
on its compact domain gives a finite subdivision
$0=r_0<\cdots<r_N=1$ such that each pair
\[
 t^{(j)}=(H(-,r_j),H(-,r_{j+1})):Y\longrightarrow C_d
\]
lands in $C_d$.  Its distinguished value has the neighborhood basis just verified.
By \Cref{prop:class-covariance,thm:homotopy},
\[
 z_i(H(-,r_j))
   =\pi_{0*}z_i(t^{(j)})
   =\pi_{1*}z_i(t^{(j)})
   =z_i(H(-,r_{j+1})).
\]
Iteration proves the assertion.
\end{proof}

\subsection{Equal parts and the fixed lattice}

Fix $k\geq1$, put $m=p^k$, and let
$\zeta=\exp(2\pi\mathrm i/m)$.  Choose a transversal $D_0$ for the
cosets of $\mu_m$ in $\Lambda$, and set
\begin{equation}\label{char:cyclic-partition}
 D_i=\zeta^iD_0\qquad(0\leq i<m).
\end{equation}
These are disjoint and partition $\Lambda$.
By \Cref{lem:finite-sheets} there are a nonempty open
$V\subset Y\setminus\{y_\infty\}$ and a locally constant function
$\beta:q^{-1}V\to S^1$ satisfying
\[
 \beta(gx)=\chi(g)\beta(x),\qquad \chi(1)=\zeta,
\]
for the character $\chi:G\to\mu_m$.

\begin{proposition}[Equality for a concentrated test]\label{prop:equal-parts}
Let $k\geq1$, and let $V$ and its character multiplier $\beta$ be as
in \Cref{lem:finite-sheets}. If $h:Y\to S^d$ is constant on
$Y\setminus V$, then the classes for
\eqref{char:cyclic-partition} satisfy
\[
 z_0(h)=z_1(h)=\cdots=z_{m-1}(h)\quad\text{in }E_d(S^d).
\]
\end{proposition}

\begin{proof}
Let $b$ be that constant value.  Since $y_\infty\notin V$, it equals
$b_h$.  The source of the restricted test over $S^d\setminus\{b\}$
lies in $q^{-1}V\subset W$.  Multiplication by $\beta$ on its
coefficient sheaf is defined because $\beta$ is locally constant, and
is a real isometry because $|\beta|=1$.  With our convention for $T$,
\[
 T M_\beta=M_{\zeta^{-1}}M_\beta T.
\]
It therefore takes a $\lambda$-eigensheaf to the
$\zeta^{-1}\lambda$-eigensheaf and cyclically permutes the subsums
in \eqref{char:cyclic-partition}.  After the proper restricted
manifold pushforward it remains an isometry of
$(A_h,\alpha_h)|_{S^d\setminus\{b\}}$ and conjugates $e_i$ to
$e_{i-1}$, with indices modulo $m$.

The isometry and its inverse define quotient morphisms by
\Cref{thm:germs}.  Its conjugation identity makes the twisted forms
for $e_i$ and $e_{i-1}$ isometric in the quotient.  Their untwisted
terms agree as well.  The integral isomorphism in
\Cref{prop:contractible-quotient} and
\Cref{def:character-classes} now give $z_i(h)=z_{i-1}(h)$.
\end{proof}

\begin{theorem}[Contradiction on the invariant chart]\label{thm:contradiction}
For every $n\geq1$ and every prime $p$, an orientation-preserving
continuous $\Zp$-action on a connected open subset of $\R^n$ has
nonzero kernel.
\end{theorem}

\begin{proof}
Suppose such an action is faithful.  Choose an odd integer $d>n$ with
$d\geq3$, and form the data at the beginning of this section by
\Cref{lem:collapse,lem:averaged-test}.  The latter may replace $G$ by
an open subgroup; its action remains faithful and orientation preserving.

First fix the lattice.  \Cref{thm:sphere-lattice} gives a positive
power of $2$, denoted by $L_d$, such that
\[
 E_d(S^d)\otimes\Q=\Q\overline{u}_d,\qquad
 \im\bigl(E_d(S^d)\longrightarrow E_d(S^d)\otimes\Q\bigr)
       \subset L_d^{-1}\Z\overline{u}_d.
\]
This bound depends only on the chosen sphere comparison.  Choose $k$
so large that $m=p^k>4L_d$, and then choose $V,\beta$ and the partition
\eqref{char:cyclic-partition}.  These choices give integral classes
$z_i(a)\in E_d(S^d)$.

Write $\overline{z}_i(a)=z_i(a)\otimes1$ and
$\overline{u}_d=u_d\otimes1$.  We now compare these rational images. By \Cref{prop:move-test} there
are an integer $s>0$ and a map $h$,
constant on $Y\setminus V$, such that $h\simeq D_s a$, where
$D_s:S^d\to S^d$ has degree $s$.  Thus
\Cref{prop:equal-parts,prop:class-homotopy,prop:class-covariance} give
\[
 D_{s*}z_i(a)=z_i(D_s a)=z_i(h)=z_j(h)=D_{s*}z_j(a).
\]
On the rational line $\Q\overline{u}_d$, the map $D_{s*}$ is multiplication by
$s$, by \Cref{prop:degree-action}.  Since $s>0$, all the $\overline{z}_i(a)$
are equal.  This cancels $s$ only as a nonzero
rational scalar: the original classes $z_i(a)$ remain integral, and
their rational images lie in the same fixed lattice.  In particular,
the bound $L_d$ is unchanged,
regardless of the size or prime divisors of $s$.

Finally compute their sum.  The coefficient form $B$ is the orthogonal
sum of two positive real lines.  Therefore
$[A_a,\alpha_a]=2(aqc)_*u_d$, and the degree-one property of $aqc$
gives $\overline{[A_a,\alpha_a]}=2\overline{u}_d$.  By
\Cref{prop:sum-classes},
\[
 \sum_{i=0}^{m-1}\overline{z}_i(a)=4\overline{u}_d,
 \qquad \overline{z}_i(a)=\frac4m\overline{u}_d
       \quad\text{in }E_d(S^d)\otimes\Q.
\]
Each $z_i(a)$ is an integral class, so
$\overline{z}_i(a)\in L_d^{-1}\Z\overline{u}_d$. The lattice inclusion
would force $4L_d/m\in\Z$.  But $0<4L_d/m<1$.  This contradiction applies
to every prime, including $p=2$.
\end{proof}

\begin{proof}[Proof of \Cref{thm:main}]
Let $M$ have any positive finite dimension $n$, and suppose that its
given action were faithful.  \Cref{prop:chart-reduction} produces a faithful
orientation-preserving action of an open subgroup, itself isomorphic
to $\Zp$, on a connected invariant chart subset $O\subset\R^n$.
This contradicts \Cref{thm:contradiction}.  Hence the original kernel
is nonzero.  It is closed, since it is the intersection of the closed
stabilizers of the points of $M$.  Every nonzero closed subgroup of
$\Zp$ is open: if it contains $x$ of valuation $j$, it contains
$\overline{x\Z}=p^j\Zp$, the closure of the integer multiples of $x$.
The action therefore factors through a finite quotient.  The chart reduction
requires none of compactness, global orientability, or triangulability
of $M$, and permits all stabilizers, as established in
\Cref{sec:local-tests}.
\end{proof}

\begin{proof}[Proof of \Cref{cor:lie}]
The classical Hilbert--Smith reduction says that a locally compact
group which acts faithfully and continuously on a connected manifold
and is not a Lie group contains a topologically embedded copy of
$\Zp$ for some prime $p$; see \cite[Introduction, p.~1]{Pardon2013}.
Apply this reduction to the given locally compact Hausdorff
second-countable group.  Restriction to the embedded subgroup retains
joint continuity and faithfulness on the same $n$-manifold, contrary
to \Cref{thm:main}.  The group is therefore a Lie group.
\end{proof}
